\section{The original coefficient space and the exact minor weights}
\label{sec:coefficient-fock-bridge}

\status{CFB1--CFB13 reconstruct the two finite coefficient interfaces
used by the simultaneous construction of the original
\texttt{JointHarm.\allowbreak exists\_ioP}.  They prove the literal
\texttt{CoefReal.\allowbreak exists\_poly\_model}, the unitary cyclicity of its
original extremal polynomial, the complete arithmetic
\texttt{ArithAt}, and \texttt{weights\_of\_minors}.
The maps are built from the supplied \texttt{c0L}, \texttt{cV} and
\texttt{BV} of one original datum.  Their coefficient image is exactly
\(\mathrm{VV}_{\mathbb C}(\lambda)\), and their covariance uses the
original real action.  All three integer character exponents, both
signed weight identities and the original choice of
\texttt{lamLdef} are identified.  These are conditional implications
on the original family and model predicates; this chapter does not
construct that common datum or assert the analytic theta extension.}

All indices below start at zero, as in the source, except that
\texttt{IndefiniteWeightData.WeightData} explicitly uses one-based
arguments.  The source pin is
\begin{center}
\texttt{7cbf7a0c32fca1badf05931879a37c3f8a2b03b4}.
\end{center}
The formal proof is assumed correct.  Each displayed source contract
is accompanied by its human argument, rather than used as a proof of
itself.

\subsection*{The literal polynomial representation}

Let \(I\) be a finite index set, let \(n\ge0\), and let
\(\lambda_\tau=(\lambda_{\tau,0},\ldots,
\lambda_{\tau,n-1})\) be integral dominant weights.
Put \(d_\tau=\lambda_\tau.\mathrm{detShift}\) and
\(\nu_\tau=\lambda_\tau.\mathrm{natW}\).  Thus for \(n>0\),
\(d_\tau=\lambda_{\tau,n-1}\) and
\(\nu_{\tau,a}=\lambda_{\tau,a}-d_\tau\); at rank zero use the
source's zero shift and constant polynomial.  In separate square
matrix variables \(X_\tau\), the space
\(\mathrm{VV}_K(\lambda)\) is the image of the tensor product of the
literal \(\mathrm{nablaSq}_K(n,\lambda_\tau)\) spaces under
multiplication.  HP1 identifies it with the corresponding finite
polynomial kernel, including its rational structure.

For a tuple of invertible matrices write
\[
 \begin{split}
 (\rho_K(g)F)(X)&=
       \left(\prod_{\tau\in I}\det(g_\tau)^{d_\tau}\right)
             F((g_\tau^{\mathsf t}X_\tau)_\tau),\\
 (\mathrm{vact}_\lambda(u)F)(X)&=
       \left(\prod_{\tau\in I}\det(u_\tau)^{-d_\tau}\right)
             F((u_\tau^{-1}X_\tau)_\tau),\\
 T(g)_\tau&=(g_\tau^{-1})^{\mathsf t}.
 \end{split}                                                    \tag{CFB.1}
\]
These are exactly \texttt{rhoK}, \texttt{vact} and
\texttt{invTrK}.  The first two expressions act on the full polynomial
ring, and \(\rho_K\) restricts to \(\mathrm{VV}_K\).
All determinant powers are integer powers, including negative ones.

% BEGIN CFB1
\paragraph{CFB1: base change and the exact inverse-transpose action.}
\label{claim:CFB1}
For every field embedding \(\phi:K\hookrightarrow\mathbb C\),
coefficientwise base change induces the source semilinear map
\[
 b_\phi=\mathrm{mapVV}(\phi,\lambda):
        \mathrm{VV}_K(\lambda)\longrightarrow
        \mathrm{VV}_{\mathbb C}(\lambda).
                                                               \tag{CFB.2}
\]
The images under \(b_\phi\) of the original basis
\(\mathrm{bVV}_K(\lambda)\) form a complex basis, and
\[
 b_\phi(\rho_K(g)v)=\rho_{\mathbb C}(\phi(g))b_\phi(v),
 \qquad
 \mathrm{vact}_\lambda(u)=\mathrm{actP}_\lambda(T(u)).
                                                               \tag{CFB.3}
\]
The latter identity holds on the full polynomial ring and hence on
the unchanged subspace \(\mathrm{VV}_{\mathbb C}(\lambda)\).

\emph{Proof.}
The embedding into \(\mathbb C\) makes \(K\) characteristic zero.
By HP1, each factor and then the tensor polynomial space is the
solution of finite linear equations with rational coefficients in
its prescribed finite multidegree space.  Gaussian elimination over
\(\mathbb Q\) commutes with extension to \(K\) and then to
\(\mathbb C\).  Thus coefficientwise base change of any \(K\)-basis
is a complex basis of the same literal polynomial kernel.
This applies in particular to \(\mathrm{bVV}_K\), without replacing
it by the row-reduction basis.  Equivalently, a nonzero maximal
coefficient minor of that basis stays nonzero under \(\phi\), and
the kernel dimensions agree by the same elimination.  The separate
variable blocks make the tensor multiplication map injective, so
the conclusion concerns its exact polynomial image.

Coefficientwise application of \(\phi\) commutes with finite sums,
products, determinants, inverses of nonzero determinants and matrix
substitution.  This proves the first identity in (CFB.3).
For the second, transpose twice in (CFB.1) and use
\(\det T(u_\tau)=\det(u_\tau)^{-1}\).
Finally \(T(gh)=T(g)T(h)\): inverse and transpose each reverse
order, and their composition preserves it.  Consequently both
actions in (CFB.1) are group representations on the stated space.
These calculations also prove that every coefficient of
\(\mathrm{actP}_\lambda(g)F\), for fixed \(F\), is a finite sum
of products of matrix entries and integer powers of nonzero
determinants.  It is continuous for a continuous tuple of invertible
matrices.  The arguments include \(n=0\) and an empty \(I\), where
the polynomial space is the one-dimensional constant space.
% END CFB1

% BEGIN CFB2
\paragraph{CFB2: the complex extension determined by the supplied coefficient map.}
\label{claim:CFB2}
Let \(W\) be a complex vector space and let
\(c_0:\mathrm{VV}_K(\lambda)\to W\) be \(\phi\)-semilinear.
Assume that the vectors
\(c_0(\mathrm{bVV}_K(\lambda)_a)\) are complex-linearly independent
and span \(W\).  There is a unique complex linear isomorphism
\[
 e:\mathrm{VV}_{\mathbb C}(\lambda)\xrightarrow{\ \sim\ }W,
 \qquad e(b_\phi(v))=c_0(v).                                  \tag{CFB.4}
\]
In particular this is the exact \texttt{exists\_scalarExt} interface.

\emph{Proof.}
Write \(b_a=\mathrm{bVV}_K(\lambda)_a\).  CFB1 makes
\((b_\phi(b_a))_a\) a complex basis.  Define
\[
 e\left(\sum_a z_a b_\phi(b_a)\right)=\sum_a z_a c_0(b_a).
                                                               \tag{CFB.5}
\]
The assumed independence and spanning give its inverse by the same
coordinate formula.  For \(v=\sum_a k_a b_a\), semilinearity gives
both sides of (CFB.4) as \(\sum_a\phi(k_a)c_0(b_a)\).
Any other complex linear map with (CFB.4) agrees on a basis, so is
equal to \(e\).  This is also the elementary extension mechanism of
SE1, applied to the specified basis.  There is no density assumption
on the coefficient field and no choice of a scalar or an invariant
inner product in this construction.
% END CFB2

% BEGIN CFB3
\paragraph{CFB3: rational points with the original finite part removed.}
\label{claim:CFB3}
Let \(V\) be the original Hermitian space over the original CM
extension \(E/F\), and retain \(\theta:\mathrm{InfFam}(D)\)
with its \(\mathrm{IsInfFam}\) condition.  For every
\(u\in G_\infty(V)=\ker(\mathrm{finHom})\) and every open
neighborhood \(\mathcal U\) of \(u\) in the adelic unitary group,
there exist \(\gamma\in G_{\mathbb Q}(V)\) and
\(k\in G_f(V)=\ker(\mathrm{infHom})\) such that
\[
             \gamma k\in\mathcal U\cap G_\infty(V).             \tag{CFB.6}
\]
These are the literal subgroups and topology of
\texttt{weak\_approx\_real}.

\emph{Proof.}
HCU14 identifies the original infinite subgroup continuously, with
continuous inverse, with the product of the unitary groups of the
original \(\mathrm{gramAt}(D,V,\tau)\).  Its proof identifies the
given \(\theta\) entry by entry; it does not choose replacement
embeddings.  Rational unitary matrices are dense in this finite
product by the Cayley proof in MF3.  To recall why the density uses
only weak approximation, the skew-adjoint matrices for the fixed
Hermitian form are an \(F\)-vector space.  The rational map
\(Z\mapsto(1-Z)(1+Z)^{-1}\) is an isomorphism between its open
chart and the chart \(\det(1+g)\ne0\) of the unitary group.
At each real place a norm-one scalar arbitrarily close to \(1\)
avoids the finitely many values for which the latter determinant
vanishes.  The one-dimensional norm-one Cayley parametrization and
weak approximation choose these scalars simultaneously.  Weak
approximation in the skew-adjoint vector coordinates then gives
rational matrices arbitrarily close to all the prescribed local
matrices.  The coordinate approximation itself follows from density
of \(F\) in its real embedding product: the embedding matrix of a
rational basis is invertible by separability, and
\(\mathbb Q^{[F:\mathbb Q]}\) is dense in its real vector space.

Choose such a rational \(\gamma\) sufficiently close to the
archimedean tuple of \(u\).  In the product of finite and infinite
adelic rings, take \(k\) to have finite component
\(\gamma_f^{-1}\) and infinite component \(1\).  It satisfies
the same unitary equation componentwise, so belongs to the original
\(G_f(V)\).  The product \(\gamma k\) has tuple
\((\gamma_\infty,1_f)\).  The continuous inverse in HCU14 makes
this tuple belong to \(\mathcal U\) once the archimedean
approximation is sufficiently close.  Its finite component is
exactly the identity, proving (CFB.6).  Nothing here requires the
rational points themselves to approximate the finite component.
% END CFB3

% BEGIN CFB4
\paragraph{CFB4: the same coefficient map intertwines every original real point.}
\label{claim:CFB4}
Fix the original \(D,V,\theta,c,\lambda,c_0\), where
\(c:\mathrm{CoefA}(V)\), and assume
\[
 \mathrm{IsInfFam}(\theta),\qquad c.\mathrm{Is},\qquad
 \mathrm{CoefMapTie}\bigl(V,c.\mathrm{coef},c_0,
                  \rho_{Vr}(\mathrm{labLL}(D),\lambda)\bigr).
                                                               \tag{CFB.7}
\]
Here \(K=D.\mathrm{what}_2.\mathrm{LL}\) and
\(\phi=D.\mathrm{what}_2.\mathrm{embLL}\).
The unique \(e\) of CFB2 satisfies, for every original
\(u\in G_\infty(V)\) and \(v\in\mathrm{VV}_{\mathbb C}(\lambda)\),
\[
       c.\rho(u)e(v)
       =e\bigl(\rho_{\mathbb C}(T(R(u)))v\bigr),\qquad
 R(u)_\tau=\mathrm{realGL}(\theta,V,u)_\tau.
                                                               \tag{CFB.8}
\]
This proves the whole \texttt{exists\_real\_model} output, with
the additional uniqueness of its coefficient extension.

\emph{Proof.}
The last two clauses of the given \texttt{CoefMapTie} are precisely
the basis independence and spanning in CFB2.  For a rational
\(\gamma\in G_{\mathbb Q}(V)\), take its actual matrix
\(\gamma_E\) satisfying \texttt{IsRatMat}; existence is part of
the definition of that rational subgroup.  The original field tie
gives, for every \(\tau\),
\[
        \phi\circ\mathrm{labLL}(D,\tau)
                           =\mathrm{tauE}(D,\tau).             \tag{CFB.9}
\]
Indeed both sides evaluate \(x\) as
\(\phi_0(\mathrm{sigmaOf}(D.eTau(\tau))(x))\), by
\(D.\mathrm{embE}_{\phi_0}\).
The \(\mathrm{IsInfFam}\) condition says that evaluation of the
diagonal rational matrix by \(\theta_\tau\) is this same
\(\mathrm{tauE}(D,\tau)(\gamma_E)\).
Base change commutes with inverse and transpose, so CFB1 identifies
the base change of the exact source action
\(\rho_{Vr}(\mathrm{labLL},\lambda)(\gamma_E)\) with
\(\rho_{\mathbb C}(T(R(\gamma)))\).
The rational covariance clause of the given tie and (CFB.4)
therefore give (CFB.8) for \(v\) in the range of \(b_\phi\).
That range spans by CFB1, so complex linearity gives the identity
for every \(v\).

If \(k\in G_f(V)\), the first clause of \(c.\mathrm{Is}\)
gives \(c.\rho(k)=1\), while \(R(k)=1\) by definition of
\(G_f(V)\).  Thus (CFB.8) also holds for every \(\gamma k\)
as in CFB3.  For fixed \(v\), both sides of (CFB.8) are continuous
functions of the adelic group element: the left side uses the orbit
continuity explicitly supplied by \(c.\mathrm{Is}\), and the right
side uses HCU14 and the finite coefficient continuity in CFB1.
Equivalently, apply every polynomial coefficient functional to
\(e^{-1}\) of the left side; all these functionals are continuous
on the finite-dimensional coefficient space.  CFB3 makes the points
\(\gamma k\) dense at each prescribed \(u\in G_\infty(V)\).
Equality of the continuous coordinate functions there proves
(CFB.8).  Irreducibility and an arbitrary real intertwiner have not
been introduced as extra premises.
% END CFB4

% BEGIN CFB5
\paragraph{CFB5: the original-frame polynomial map, with exact image.}
\label{claim:CFB5}
Under (CFB.7), let \(B=(B_\tau)_\tau\) be any specified tuple of
invertible complex matrices of the original rank.  There is an
explicit injective complex linear map
\(j:c.W\to\mathrm{Pol}_{\mathbb C}(I,n)\) such that
\[
 \begin{split}
 \operatorname{range}(j)&=\mathrm{VV}_{\mathbb C}(\lambda),\\
 j(c.\rho(u)w)&=
   \mathrm{vact}_\lambda\bigl((B_\tau^{-1}
          \mathrm{compAt}(\theta,\tau,\mathrm{mt}(V,u))
                    B_\tau)_\tau\bigr)j(w)
       \quad(u\in G_\infty(V)),\\
 j(c_0(v))&=\mathrm{incl}\,
       \rho_{\mathbb C}(T(B^{-1}))b_\phi(v).
 \end{split}                                                    \tag{CFB.10}
\]
In particular the first two lines are the literal
\texttt{CoefReal.exists\_poly\_model}.

\emph{Proof.}
Use the \(e\) determined in CFB2--CFB4 and define
\[
          j=\mathrm{incl}\circ
               \rho_{\mathbb C}(T(B^{-1}))\circ e^{-1}.
                                                               \tag{CFB.11}
\]
The middle factor has inverse \(\rho_{\mathbb C}(T(B))\), by
the group law verified in CFB1.  Hence it is an automorphism of
the unchanged \(\mathrm{VV}_{\mathbb C}(\lambda)\).
The composite before inclusion is an isomorphism onto that very
space; both injectivity and the exact range in (CFB.10) follow.
The last line follows by applying (CFB.4).

For covariance, abbreviate \(R=R(u)\).  The matrix of \(R\)
is exactly \(\mathrm{compAt}(\theta,\tau,\mathrm{mt}(V,u))\)
in each component, by its definition.  CFB4 makes the two operator
products in the desired equality, before inclusion,
\[
 \rho(T(B^{-1}))\rho(T(R))e^{-1},\qquad
 \rho(T(B^{-1}RB))\rho(T(B^{-1}))e^{-1}.
                                                               \tag{CFB.12}
\]
Both equal \(\rho(T(B^{-1}R))e^{-1}\); use CFB1 to identify
the second expression with \(\mathrm{vact}\) on the polynomial
image.  This proves (CFB.10).  No isometry of \(j\) is asserted:
the subsequent unit-vector normalization belongs to the Fock-map
construction.  The present map is fixed by the supplied \(c_0\)
and the specified \(B\).
% END CFB5

% BEGIN CFB6
\paragraph{CFB6: the unchanged extremal polynomial is cyclic under the compact source.}
\label{claim:CFB6}
Let
\[
 h_\lambda=\mathrm{ten}\bigl(\tau\mapsto
              \mathrm{extremal}_{\mathbb C}(n,\lambda_\tau)\bigr).
\]
Then the exact compact orbit spans the exact coefficient space:
\[
 \mathrm{VV}_{\mathbb C}(\lambda)
   =\operatorname{span}_{\mathbb C}
       \{\mathrm{vact}_\lambda(u)h_\lambda:
                                  u\in\prod_{\tau\in I}U(n)\}.
                                                               \tag{CFB.13}
\]
The compact group here is the source's
\(\prod_\tau\mathrm{Kset}(\mathrm{allPos}(n))\), so this supplies
\texttt{VV\_le\_span\_\allowbreak vact\_extremal} on its literal polynomial.

\emph{Proof.}
HP1 identifies \(h_\lambda\) with the product of the leading
minor monomials in the separate blocks.  It belongs to the literal
\(\mathrm{VV}\), and its value at the identity tuple is \(1\).
The orbit span \(S\) is therefore nonzero, is contained in
\(\mathrm{VV}\), and is stable under \(\mathrm{vact}(\prod U(n))\).
We show directly that this implies stability under the full complex
general linear product.

In each block, the real skew-Hermitian space has the real basis
\[
 iE_{aa},\qquad E_{ab}-E_{ba},\qquad
                    i(E_{ab}+E_{ba})\quad(a<b).
                                                               \tag{CFB.14}
\]
The same list is a complex basis of \(M_n(\mathbb C)\): its
off-diagonal pairs recover \(E_{ab},E_{ba}\), and its diagonal
vectors recover \(E_{aa}\).  For real coordinates \(z\) in this
basis, \(C(z)=(1-Z)(1+Z)^{-1}\) is unitary.  Indeed \(Z^*=-Z\),
and \(1+Z\) is invertible because a skew-Hermitian matrix has no
real eigenvalue \(-1\), as also follows by taking the real part
of \(\langle Zv,v\rangle\).  The Cayley formula then verifies
\(C(z)^*C(z)=1\).

Choose a basis of \(S\) and any linear quotient map with kernel
\(S\).  Applying that quotient map to
\(\mathrm{vact}(C(z))v\), for a basis vector \(v\in S\), gives
finitely many rational functions of the coordinates.  All
denominators are products of powers of \(\det(1\pm Z_\tau)\),
including those needed for the integer determinant twists.
They do not vanish at real skew-Hermitian coordinates.  Each
numerator vanishes for every real coordinate tuple, by compact
stability.  Induction in the number of variables using the
one-variable polynomial identity principle makes that numerator
the zero polynomial over \(\mathbb C\).
Because (CFB.14) is also a complex basis, complex Cayley coordinates
parametrize the open subset \(\det(1+g_\tau)\ne0\) of
\(\prod\mathrm{GL}_n(\mathbb C)\), with inverse
\((1-g)(1+g)^{-1}\).  Thus every quotient coefficient vanishes
on that open set and hence on the whole general linear product:
after clearing determinants it is a polynomial vanishing on a
nonempty polynomial open subset, and a polynomial ring is a domain.
This proves full \(\mathrm{vact}\)-stability of \(S\).

The map \(T\) in CFB1 is a bijective group automorphism, so this
is equivalent to stability under all \(\rho_{\mathbb C}(g)\).
HP2, whose proof is the finite factorial-pairing and highest-line
argument on this same polynomial module, now gives
\(S=\mathrm{VV}\).  Negative determinant shifts cause no change:
their factors are nonzero regular functions on the invertible
group and have already been retained when clearing denominators.
For an empty \(I\) or rank zero, both sides of (CFB.13) are the
constant line and the assertion holds directly.
% END CFB6

\subsection*{The original floors, dominant weight, and signed lowest type}

For the rest of the finite arithmetic fix the original \(D\),
\(i\in\{0,\ldots,m-1\}\), and \(\tau\).  Write
\[
 \begin{gathered}
 k=D.R.k,\quad m=D.R.m,\quad t=D.R.t,\quad r=k+m-1,\\
 p=\mathrm{sig}(W_c(D,i),\mathrm{tauE}(D,\tau))_1,
 \quad q=\mathrm{sig}(W_c(D,i),\mathrm{tauE}(D,\tau))_2,\\
 n=\mathrm{rank}(V_c(D,i)),\quad e=(D.\mathrm{hc}_\tau).e_i,
 \quad \ell_a=\mathrm{lamAt}(D.\mathrm{hc}_\tau,a),\\
 \alpha_a=\mathrm{aPl}(D,i,p,\tau,a),\qquad
 \beta_a=\mathrm{aMi}(D,i,p,q,\tau,a).
 \end{gathered}                                                 \tag{CFB.15}
\]
The notation \(\beta\) for the negative minor degrees is local to
this section and is unrelated to a level matrix in the original
datum.  Values \(\alpha_a\) outside \(a<p\) and \(\beta_a\)
outside \(a<q\) are zero by their actual definitions.

% BEGIN CFB7
\paragraph{CFB7: the original ranks and integral decreasing row.}
\label{claim:CFB7}
The source fields give
\[
 p=k+i,\qquad q=m-i-1,\qquad p+q=r=\mathrm{rank}(W_c(D,i)),
 \qquad n=r+t=p+q+t.                                          \tag{CFB.16}
\]
Here \(r,t\) are odd, \(n\) is even, \(p>0\), and \(t>0\).
The entries \(e\) and \(\ell_a\), \(0\le a<r\), are integers,
and for \(0\le a\le b<r\),
\[
                    \ell_b+(b-a)\le\ell_a.                    \tag{CFB.17}
\]
Moreover
\[
 \ell_{p-1}=-b_i,\qquad e-b_i\ge(t+3)/2,
 \qquad
 q>0\ \Longrightarrow\ 
       \ell_p=-b_{i+1},\quad b_{i+1}-e\ge(t+3)/2,
                                                               \tag{CFB.18}
\]
where \(b_j=(D.\mathrm{hc}_\tau).b_j\).

\emph{Proof.}
Insert \(j=i+1\) into the original \(W\)-signature and
\(W,V\)-rank fields.  Since \(0\le i<m\), all their stated
index guards hold; ordinary subtraction gives (CFB.16).
The fields \(k,m\) even and \(t\) odd make \(r=k+m-1\) odd
and \(n=r+t\) even.  The rank fields include \(m\ge2\),
\(k\ge m+3\) and \(t>3m\), which give the displayed positivity.

The actual row is
\[
       (v_0,\ldots,v_{k-2},-b_0,\ldots,-b_{m-1}).              \tag{CFB.19}
\]
The \(\mathrm{HCParam}\) fields make each entry integral.
Successive \(v\)'s satisfy
\(v_{a+1}<u_{a+1}<v_a\).
At the join they satisfy \(-b_0<u_{k-1}<v_{k-2}\).
For successive tail entries the two gap fields give
\(b_j<e_j<b_{j+1}\), hence \(-b_{j+1}<-b_j\).
Thus (CFB.19) is strictly decreasing.  Strict decrease of integers
means a drop of at least one per step; summing these drops proves
(CFB.17).  The last positive-block index is
\(p-1=(k-1)+i\), giving its exact tail value in (CFB.18).
If \(q>0\), then \(i+1<m\), so the next tail value and the
second gap are defined and give the remaining identities.
% END CFB7

% BEGIN CFB8
\paragraph{CFB8: every floor disappears, with both degree sequences antitone.}
\label{claim:CFB8}
Define the rational expressions
\[
 \begin{aligned}
 A_a&=\ell_a+e-(p-a-1)-(t+1)/2 &&(0\le a<p),\\
 D_c&=\ell_{p+c}+e+c+1+(t-1)/2 &&(0\le c<q).
 \end{aligned}                                                  \tag{CFB.20}
\]
They are integers, and the literal source degrees satisfy
\[
       \alpha_a=A_a\ge1,\qquad
       \beta_{q-1-c}=-D_c\ge1.                                \tag{CFB.21}
\]
Both \(\alpha\) on \(0,\ldots,p-1\) and \(\beta\) on
\(0,\ldots,q-1\) are antitone.  In particular all the positivity,
floor, and antitonicity fields in the original \texttt{ArithAt}
hold without a further weight-arithmetic premise.

\emph{Proof.}
Since \(t\) is odd, \((t+1)/2\) and \((t-1)/2\) are integers.
Together with CFB7 this gives integrality of (CFB.20).
For \(a\le b<p\), (CFB.17) gives
\[
 A_b-A_a=(\ell_b+b)-(\ell_a+a)\le0.
\]
Likewise \(D_d\le D_c\) for \(c\le d<q\).
At the two ends, (CFB.18) gives
\[
 A_{p-1}=e-b_i-(t+1)/2\ge1,
 \qquad D_0=e-b_{i+1}+(t+1)/2\le-1\quad(q>0).
                                                               \tag{CFB.22}
\]
Thus every \(A_a\ge1\), and every defined \(-D_c\ge1\).

The actual \(\alpha_a\) is
\(\max(\lfloor A_a\rfloor,0)\) on \(a<p\).
For an integral nonnegative rational number this equals the number
itself, proving the first identity in (CFB.21).
For the negative degrees, substitute \(a=q-1-c\) into the actual
definition of \(\mathrm{aMi}\).  Natural-number subtraction is
exact because \(0\le c<q\): \(q-a=c+1\), and
\(p+(q-a)-1=p+c\).  The expression inside its floor becomes
\[
          -[e+\ell_{p+c}+(c+1)+(t-1)/2]=-D_c.
\]
It too is a positive integer, proving the second identity.
As \(a\) increases, \(q-1-a\) decreases; antitonicity of
\(D_c\) therefore makes \(-D_{q-1-a}=\beta_a\) antitone.
If \(q=0\), every negative-block assertion is empty and the
definition gives \(\beta_a=0\) for all \(a\).
% END CFB8

% BEGIN CFB9
\paragraph{CFB9: the dominant minor weight forces the original stored choice.}
\label{claim:CFB9}
Put
\[
 d=(p-q)/2-e-1/2=(p-q-1)/2-e,
 \qquad
 \Theta_a=
 \begin{cases}
       \alpha_a,&0\le a<p,\\
       0,&p\le a<p+t,\\
       -\beta_{q-1-c},&a=p+t+c,\quad0\le c<q.
 \end{cases}                                                    \tag{CFB.23}
\]
Then \(d\in\mathbb Z\), and \(L_a=\Theta_a+d\), \(0\le a<n\),
is an integral dominant weight.  For the \emph{actual} source choice
\[
        \lambda=\mathrm{Bridge.lamLdef}(n,i+1,\tau)
\]
one has
\[
 \mathrm{IsLamPrime}(\lambda,i+1,\tau),\qquad
                  \lambda.w_a=\Theta_a+d\quad(0\le a<n).
                                                               \tag{CFB.24}
\]
In particular the fallback branch of \texttt{lamLdef} cannot occur.

\emph{Proof.}
Since \(p+q\) is odd, \(p-q-1=(p+q)-2q-1\) is even.
CFB7 therefore makes \(d\) an integer.  CFB8 makes the first
block of \(\Theta\) antitone with values at least one.  Its
middle block is zero.  In the last block,
\(-\beta_{q-1-c}=D_c\le-1\) is antitone as \(c\) increases.
These statements check every internal and boundary comparison,
including \(q=0\), and prove dominance of \(\Theta\) and \(L\).

It remains to identify the original conditional choice, rather
than merely exhibit a different dominant weight.  Choose integral
representatives of the supplied \(e,v_a,b_a\), and define a
\texttt{WeightData} \(W\) with exactly
\[
 W.k=k,\ W.m=m,\ W.t=t,\ W.j=i+1,\ W.e=e,\quad
 W.v(a+2)=v_a,\quad W.b(a+1)=b_a.                             \tag{CFB.25}
\]
Give unused natural indices arbitrary integral values, for example
zero.  They do not enter any formula below.  Then
\(W.p=p\), \(W.q=q\), \(W.nV=n\), and for \(0\le a<r\),
\[
                         W.\mathrm{lam}(a+1)=\ell_a.           \tag{CFB.26}
\]
For \(a<k-1\) this is \(W.v(a+2)=v_a\); for
\(a=(k-1)+j\), it is \(-W.b(j+1)=-b_j\).

We verify all three blocks of the one-based source formula.
For \(a<p\), substitution into \(W.\mathrm{aPlus}\) gives
\[
 \begin{split}
 W.\mathrm{aPlus}(a+1)
 &=\ell_a+e-\frac{n+1-2(a+1)}2-\frac{p-q}2\\
 &=\ell_a+e-(p-a-1)-(t+1)/2=A_a=\alpha_a.
 \end{split}                                                    \tag{CFB.27}
\]
For \(a=p+t+c\), \(0\le c<q\), the argument of
\(W.\mathrm{aMinus}\) in \(W.\mathrm{thetaWt}(a+1)\) is
\(s=q-c\).  Consequently \(q+1-s=c+1\), and
\[
 \begin{split}
 W.\mathrm{aMinus}(q-c)
 &=b_{i+c+1}-e+\frac{1-t-2(c+1)}2\\
 &=b_{i+c+1}-e-c-(t+1)/2=-D_c=\beta_{q-1-c}.
 \end{split}                                                    \tag{CFB.28}
\]
Here \(i+c+1<m\) follows from \(c<q=m-i-1\), and
\(\ell_{p+c}=-b_{i+c+1}\) follows from (CFB.19).
The middle block of \(W.\mathrm{thetaWt}\) is zero by definition,
and \(W.\mathrm{dTau}=d\).  Thus, throughout the actual rank,
\[
       W.\mathrm{thetaWt}(a+1)=\Theta_a,
       \qquad W.\mathrm{lamPrime}(a+1)=\Theta_a+d.              \tag{CFB.29}
\]
The integral dominant weight \(L\) and this very \(W\) witness
\(\mathrm{IsLamPrime}(L,i+1,\tau)\).  The existential test in
\texttt{lamLdef} is therefore true.

For completeness, (CFB.26)--(CFB.29) use only the fields in
(CFB.25).  Hence they hold for \emph{every} \(W\) allowed by
\(\mathrm{IsLamPrime}\), regardless of its unused entries.
Any dominant weight satisfying that predicate has entries exactly
\(\Theta_a+d\), and so equals \(L\).
The source's selected witness consequently has (CFB.24).
This proves both existence and independence of the internal choice,
and also the universal \texttt{lamPrime\_eq} field of
\texttt{ArithAt}.
% END CFB9

% BEGIN CFB10
\paragraph{CFB10: the exact signed enumeration and lowest compact weight.}
\label{claim:CFB10}
Let \(\sigma=\mathrm{embLL}\), and let \(M\) be any of the
original \(\mathrm{DSModel}(W_c(D,i),\mathrm{parAt}(W_c(D,i),\sigma))\)
models with signs ordered as in the source,
\(M.\epsilon_\tau(a)=\mathbf1_{a<p}\).
Put \(\Lambda_a=\mathrm{lowK}(M.\epsilon_\tau,M.\mathrm{lam}_\tau)(a)\).
Its stored parameters and this literal lowest weight satisfy
\[
 \begin{gathered}
       M.\mathrm{lam}_\tau(a)=\ell_a\quad(0\le a<r),\\
 \Lambda_a=\begin{cases}
       \ell_a+q/2-(p-1)/2+a,&0\le a<p,\\
       \ell_{p+c}-p/2-(q-1)/2+c,&a=p+c,\quad0\le c<q.
 \end{cases}
 \end{gathered}                                                 \tag{CFB.30}
\]
The last two quantities are the integral source values
\(\kappa^+_a=\mathrm{kapW}(\mathrm{inl}(a))\) and
\(\kappa^-_c=\mathrm{kapW}(\mathrm{inr}(c))\), and
\[
        \kappa^+_a=\alpha_a+n/2-e,\qquad
        \kappa^-_c=-\beta_{q-1-c}-n/2-e.                        \tag{CFB.31}
\]

\emph{Proof.}
We first check the base-place label, before using an enumeration.
Equation (CFB.9) gives
\(\mathrm{labE}(\mathrm{embLL},\tau)=\mathrm{tauE}(D,\tau)\).
The \(\mathrm{tauE}\)'s are distinct: injectivity of the original
\(\phi_0\), the automorphism parametrization \(\mathrm{sigmaOf}\)
and \(D.eTau\) proves this from their definitions.
No two are conjugates.  One direct check uses the original
\(\xi_0\), whose image has positive imaginary part at every
chosen \(\mathrm{tauE}\); complex conjugation reverses that sign.
Thus the index chosen by \(\mathrm{lab}(\mathrm{embLL},\rho)\)
must be \(\rho\), and only the un-conjugated branch of
\(\mathrm{parAt}\) occurs.  This is also the base-label calculation
in HCU20; it uses only \(D\) and therefore applies to the present
arbitrary \(M\).

By definition \(\mathrm{kapQ}\) is the pair of rational
expressions in (CFB.30).  They are integers: \(\ell\) is
integral, and \((q-p+1)/2\) and \(-(p+q-1)/2\) are integers
since \(p+q\) is odd.  Its floor \(\mathrm{kapW}\) therefore
equals it.  Inserting these two expressions into the actual
\(\mathrm{holEnt}\) formulas cancels their half-sums and gives
respectively \(\ell_a\) and \(\ell_{p+c}\).
The signed \(\mathrm{holPar}\) is consequently the finite set
\[
       \{(\ell_a,\mathrm{true}):0\le a<p\}
       \ \cup\ 
       \{(\ell_{p+c},\mathrm{false}):0\le c<q\}.              \tag{CFB.32}
\]
Its first coordinates are all distinct by CFB7 and are already in
strict decreasing order.  The stored \(M.\mathrm{enum}\) says
that \(M.\mathrm{lam}\) is a strict decreasing enumeration of
this same signed set.  Such an enumeration is unique: its first
element is the unique maximum first coordinate, and deleting that
element proves the statement successively for every index.
Thus its parameter at position \(a\) is exactly \(\ell_a\),
with the stated sign.  This agrees with the retained ordering field;
no permutation of the original model is performed.

The definition of \(\rho_C\) counts equal-sign indices above an
index minus those below, divided by two; \(\rho_N\) does the
same for opposite signs.  At \(a<p\), these counts give
\(\rho_C=(p-1)/2-a\) and \(\rho_N=q/2\).
At \(p+c\), they give \(\rho_C=(q-1)/2-c\) and
\(\rho_N=-p/2\).
Since \(\mathrm{lowK}=\mathrm{lam}+\rho_N-\rho_C\), this proves
all of (CFB.30), also when the negative block is empty.
Finally substitute (CFB.20)--(CFB.21) and
\(n=p+q+t\); cancelling terms gives (CFB.31).
This proves the \texttt{kap\_plus} and \texttt{kap\_minus}
fields with the exact source normalization.
% END CFB10

% BEGIN CFB11
\paragraph{CFB11: all three integer twists and both literal weight identities.}
\label{claim:CFB11}
For each original \(i\) there are integer functions
\(x_1,x_2,x_3\) on the original \(\mathrm{Tau}\) such that,
for every \(\tau\),
\[
 \begin{split}
       x_{1,\tau}&=-e_\tau+n/2,\\
       x_{2,\tau}&=-e_\tau-n/2,\\
       x_{3,\tau}&=-d_\tau
          =-\left[\frac{p_\tau-q_\tau}{2}
                              +\left(-e_\tau-\frac12\right)\right].
 \end{split}                                                    \tag{CFB.33}
\]
In (CFB.33), \(d_\tau\) is the arithmetic shift in (CFB.23),
not a newly chosen determinant parameter.  In particular the
coefficient of \(e_\tau\) in \(x_{3,\tau}\) is positive.
For every \(M\) and base embedding as in CFB10,
\[
 \begin{split}
 \mathrm{lowK}(M.\epsilon_\tau,M.\mathrm{lam}_\tau)(a)
     &=\mathrm{wtK}(p_\tau,\alpha_\tau,\beta_\tau,
                                      x_{1,\tau},x_{2,\tau})(a),\\
 \mathrm{wtU}(p_\tau,q_\tau,\alpha_\tau,\beta_\tau,x_{3,\tau})(a)
     &=-\bigl(\mathrm{Bridge.lamLdef}(n,i+1,\tau)\bigr).w_a.
 \end{split}                                                    \tag{CFB.34}
\]
The first line is read in \(\mathbb R\) using its indicated
integer casts.  Equations (CFB.33)--(CFB.34), with the original
quantifiers, prove \texttt{weights\_of\_minors}.

\emph{Proof.}
By CFB7 write \(n=2s\) for an integer \(s\).  By CFB7 and CFB9,
\(e_\tau\) and \(d_\tau\) are integral.  Define
\(x_{1,\tau}=-e_\tau+s\),
\(x_{2,\tau}=-e_\tau-s\) and \(x_{3,\tau}=-d_\tau\).
These definitions give all three rational exponent identities,
not merely their values modulo integers.

The source \(\mathrm{wtK}\), on the \(r=p+q\) target indices,
is \(\alpha_a+x_1\) for \(a<p\) and
\(-\beta_{r-1-a}+x_2\) for \(a\ge p\).
For \(a=p+c\), its reversed index is
\(r-1-a=q-1-c\), with no truncated subtraction ambiguity.
Thus (CFB.31) proves the first line of (CFB.34).

On the source indices \(0\le a<n\), the literal formula is
\[
 \mathrm{wtU}(a)=
     -\mathbf1_{a<p}\alpha_a
       +\mathbf1_{n-1-a<q}\beta_{n-1-a}-d.                     \tag{CFB.35}
\]
Its first and last supports are disjoint, because
\(n=p+t+q\) with \(t>0\).
If \(a<p\), the last indicator vanishes and (CFB.35) is
\(-\alpha_a-d=-\Theta_a-d\).
If \(p\le a<p+t\), both vanish, so it is \(-d\).
If \(a=p+t+c\), \(0\le c<q\), then
\(n-1-a=q-1-c\), so it is \(\beta_{q-1-c}-d\).
These are exactly \(-\Theta_a-d\) in every block.
CFB9 now identifies the right side with the negative of the
\emph{stored} \(\mathrm{lamLdef}\) weight, proving the second
line.  No assertion at an arbitrary conjugate coefficient embedding
has been used or is required.
% END CFB11

% BEGIN CFB12
\paragraph{CFB12: the complete arithmetic interface and actual original frames.}
\label{claim:CFB12}
For every original \(D,i,\tau\), the literal
\(\mathrm{ArithAt}(D,i,\tau)\) holds.
For one supplied original \(R:\mathrm{RealDat}(D,A,K_s)\) with
\(R.\mathrm{IsFam}\) and \(R.\mathrm{IsModels}\), choose once
the witnesses \((\mathrm{gramR},c_{0R},c_{0L},h_T)\) in its
\(\mathrm{fam\_tie}\).  Then for each \(i\), CFB5 applies to
the actual data
\[
 \begin{gathered}
 V=V_c(D,i),\quad c=R.cV(i),\quad \theta=R.\theta,\\
 \lambda_{i,\tau}=\mathrm{Bridge.lamLdef}(\mathrm{rank}(V_c(D,i)),i+1,\tau),
 \qquad c_0=c_{0L}(i),\qquad B_\tau=R.BV(i,\tau).
 \end{gathered}                                                 \tag{CFB.36}
\]
It yields one \(j_i\) with the exact image and original real action
of (CFB.10), using this same \(c_{0L}(i)\) and \(R.BV(i)\).

\emph{Proof.}
The \texttt{ArithAt} signature and rank fields are CFB7.
Its two floor identities, the two antitone clauses and the two
positive clauses are CFB8.  The two \texttt{kap} identities are
the algebraic identities (CFB.31), whose derivation from
\(\mathrm{kapQ}\) requires only CFB7--CFB8, not a choice of \(M\).
The integer shift, the universal \(W.\mathrm{lamPrime}\) identity,
the forced \(\mathrm{IsLamPrime}\), and the final coefficient
identity are CFB9.  These are all the fields of the source record.
In particular this proof does not infer the record merely from
its two weight consequences.

For (CFB.36), the \(\mathrm{inf}\) and \(\mathrm{cV}\) fields
of \(R.\mathrm{IsModels}\) supply exactly
\(\mathrm{IsInfFam}(R.\theta)\) and \((R.cV(i)).\mathrm{Is}\).
The selected \(h_T.\mathrm{mapL}(i)\) is exactly
\(\mathrm{CoefMapTie}\) for \(c_{0L}(i)\), with
\(\rho_{Vr}(\mathrm{labLL}(D),\lambda_i)\), because these
are the actual arguments of \(R.\mathrm{fam\_tie}\).
No replacement coefficient model is involved.

Finally the \(\mathrm{BV}\) field of the same model predicate gives
\[
           \mathrm{IsONBasis}(G_\tau,1,R.BV(i,\tau),n),
\]
where
\(G_\tau=\mathrm{gramAt}(D,V_c(D,i),\tau)\).
In the definition of that predicate, each of its \(n\) signed
indices is less than \(n\), so every sign is \(+1\).  Thus
\[
                R.BV(i,\tau)^*G_\tau R.BV(i,\tau)=1.           \tag{CFB.37}
\]
The matrix \(R.BV(i,\tau)^*G_\tau\) is a left inverse.
A square matrix with a left inverse over a field is invertible,
as follows by taking determinants.  Hence the exact specified
\(R.BV(i,\tau)\) is an admissible \(B_\tau\) in CFB5.
All these choices precede any parameter \(t\) of the original
family: \(c_{0L}\), \(cV\), \(BV\) and \(\lambda\) have no
such parameter.
% END CFB12

% BEGIN CFB13
\paragraph{CFB13: one simultaneous coefficient-and-weight package on the supplied datum.}
\label{claim:CFB13}
Under the same \(R.\mathrm{IsFam}\) and \(R.\mathrm{IsModels}\)
hypotheses, for every original \(i\) there are one map \(j_i\)
and three integer functions \(x_{1,i},x_{2,i},x_{3,i}\) such that
all of the following hold on the original objects, before any
choice of a family parameter:
\[
 \begin{gathered}
 \operatorname{range}(j_i)=\mathrm{VV}_{\mathbb C}(\lambda_i),
 \qquad\text{the exact real covariance (CFB.10)},\\
 p_\tau+q_\tau=\mathrm{rank}(W_c(D,i))
          \le\mathrm{rank}(V_c(D,i)),\qquad
 \alpha_\tau,\beta_\tau\text{ antitone on their own blocks},\\
 \text{all three identities (CFB.33), and both identities (CFB.34)
 for }M_i=R.MW(i,\mathrm{sg}_0),\\
 \mathrm{VV}_{\mathbb C}(\lambda_i)
      =\operatorname{span}_{\mathbb C}
          \bigl(\mathrm{vact}(\prod_\tau U(n))h_{\lambda_i}\bigr).
 \end{gathered}                                                 \tag{CFB.38}
\]
The \(j_i\) retains the last line of (CFB.10) for the one selected
\(c_{0L}(i)\), so the scalar-extension identification remains
reviewable as an equality on the original rational coefficient map.
This supplies the complete coefficient bridge, compact cyclicity,
and finite weight inputs at the call to
\texttt{exists\_io\_poly} inside \texttt{exists\_ioP}.

\emph{Proof.}
CFB12 supplies the \(j_i\)'s and the full arithmetic record;
CFB6 supplies the exact compact cyclicity.  CFB7 gives the rank
identity and inequality since \(t>0\).  CFB8 gives the required
antitonicity.  The original \(R.hv0\) identifies
\(\mathrm{sg}_0=\mathrm{sgAt}(\mathrm{embLL},A.v_0)\) with
\(\mathrm{embLL}\), as explicitly verified in HCU20.
The \(\mathrm{order}\) field of the same
\(R.\mathrm{IsModels}\) gives the original ordering of
\(M_i.\epsilon\).  Hence CFB11 applies to precisely \(M_i\)
and supplies the three integer functions and both weights.
These constructions use \(D,i,\tau,R\) and the once-selected
\(\mathrm{fam\_tie}\) witnesses.  They do not use \(t\),
the level at \(t\), or an action chosen after \(t\).
Consequently any subsequent construction depending on this
package can choose its map once for each \(i\) and then check
the desired laws for every \(t\).

The hypotheses \(R.\mathrm{IsFam}\) and
\(R.\mathrm{IsModels}\) are retained as predicates on this
\emph{one} original datum.  Their independent existence is not
inferred from a list of abstractly isomorphic coefficient spaces.
Nor does (CFB.38) alone assert \(\mathrm{IsIoW}\), its
unit-vector normalization, the given Weil-action laws, or an
analytic theta bound.  Those are separate consumers and producers
whose finite Fock interface now receives the exact supplied data.
% END CFB13

\subsection*{Closure and the remaining common-object boundary}

The new finite implications are complete: the original rational
coefficient map extends uniquely; its actual real action is
identified by rational approximation and continuity; the actual
frame gives a polynomial image equal to the literal \(\mathrm{VV}\);
and the full compact orbit of the literal extremal polynomial
spans it.  The original floor degrees, dominant coefficient choice,
signed lowest compact weight, and all determinant twists are also
determined, including the empty negative block.

The retained inputs are the same original \(D\), family and
\(R\), its \(\mathrm{IsFam}\) coefficient tie, and its
\(\mathrm{IsModels}\) actions, frames and ordering.  Nothing in
this chapter constructs an otherwise missing compatible datum.
The simultaneous joint-harmonic/Fock construction, the actual Weil
and theta-vector inputs, and the subsequent analytic regularity
remain separate tasks.  No new published theorem is required in
this component; the inherited HP polynomial proofs, MF3 weak
approximation, and HCU14/HCU20 original-object dictionaries are
used only in their already stated scopes.
